% MINT: Budget-Aware Data-Free Mixed-Precision Quantization
% via Rate-Distortion Optimization
% Target: NeurIPS 2026 / MLSys 2027

\documentclass{article}

% Use a standard ML conference style
\usepackage[final]{neurips_2024}  % NeurIPS style, swap for target venue

\usepackage[utf8]{inputenc}
\usepackage[T1]{fontenc}
\usepackage{hyperref}
\usepackage{url}
\usepackage{booktabs}
\usepackage{amsfonts}
\usepackage{amsmath}
\usepackage{amssymb}
\usepackage{nicefrac}
\usepackage{microtype}
\usepackage{graphicx}
\usepackage{xcolor}
\usepackage{algorithm}
\usepackage{algorithmic}
\usepackage{subcaption}

\title{MINT: Budget-Aware Data-Free Mixed-Precision Quantization\\for Large Language Models via Rate-Distortion Optimization}

\author{
  baa.ai
}

\begin{document}

\maketitle

\begin{abstract}
We present \textbf{MINT} (\textbf{M}emory-\textbf{I}nformed \textbf{N}-bit \textbf{T}uning), a data-free mixed-precision quantization framework that formulates per-tensor bit-width and group-size allocation as a \emph{budget-constrained optimization problem}. Given a user-specified memory budget, MINT jointly selects the (bit-width, group-size) configuration for each weight tensor by solving a Multiple-Choice Knapsack Problem (MCKP) over per-tensor rate-distortion curves. The framework introduces three key innovations: (1) \textbf{budget-targeted quantization}---users specify an exact memory target (e.g., ``fit in 4\,GB for iPhone'' or ``fit in 24\,GB for RTX 4090'') and MINT produces a near-optimal allocation for that budget, with a fitted prediction curve that estimates output quality \emph{before} running the pipeline; (2) \textbf{joint bit-width and group-size optimization} that treats group size as a first-class allocation variable, revealing that group-size selection provides larger quality gains than bit-width changes; and (3) an \textbf{SQNR safety veto} with an empirically validated 9\,dB threshold that exploits the natural gap between catastrophic 2-bit quantization (SQNR $<$ 9\,dB, PPL triples) and usable 3-bit quantization (SQNR $>$ 10\,dB). We evaluate MINT on six model families spanning 8B--109B parameters across dense and Mixture-of-Experts architectures. On Qwen3-30B-A3B, MINT achieves +0.6\% perplexity degradation versus BF16 at 19\,GB (33\% of full size) and matches prior heuristic quality at 2.6\% less storage. On Llama-4-Scout (109B MoE), MINT enables hardware-targeted deployment of the same model across memory tiers: PPL 7.703 at 58\,GB for a 64\,GB device versus PPL 7.359 at 163\,GB for a 192\,GB device, both outperforming uniform 4-bit (PPL 7.899). In matched-size comparisons against GPTQ---a calibration-based method---across three MoE families (Qwen3-30B-A3B, Qwen2-57B-A14B, Mixtral-8x7B), MINT consistently outperforms GPTQ despite being entirely data-free: $-1.7\%$, $-0.95\%$, and $-4.6\%$ perplexity respectively at identical model sizes. Scaling experiments across budget levels reveal that quality returns diminish sharply above ${\sim}$5 average bits, and that standard downstream benchmarks (ARC-Challenge, Winogrande) are insensitive to quantization quality differences---spanning only 2.7--3.4 percentage points across the 3-to-8-bit range versus 18--20\% PPL variation---confirming perplexity as the appropriate primary metric for quantization evaluation. We also demonstrate that standard mean perplexity can be misleading due to outlier sequences, and recommend reporting median perplexity and tail statistics alongside standard corpus perplexity. The entire pipeline requires no calibration data, no gradient computation, and completes in under 50 minutes on commodity hardware.
\end{abstract}

% ═══════════════════════════════════════════════════════════════════════════════
\section{Introduction}
\label{sec:intro}
% ═══════════════════════════════════════════════════════════════════════════════

Post-training quantization (PTQ) has become the primary means of deploying large language models on consumer hardware. Methods such as GPTQ~\cite{frantar2023gptq}, AWQ~\cite{lin2024awq}, and SqueezeLLM~\cite{kim2024squeezellm} achieve remarkable compression, but they share a common requirement: a representative calibration dataset. This introduces practical concerns---calibration data may be unavailable for proprietary models, the chosen distribution may not generalize to deployment domains, and calibration demands substantial compute.

Existing data-free approaches~\cite{tang2024easyquant,zhang2025mxq,badri2024hqq} typically apply uniform bit-widths or rely on single sensitivity metrics with hand-tuned thresholds. These approaches face two fundamental limitations. First, \emph{threshold-based allocation} produces fixed bit-width decisions regardless of the deployment memory budget---the user cannot specify ``quantize this model to fit in 6\,GB'' and receive a budget-optimal allocation. Second, \emph{single-point error proxies} create circularity: using 4-bit reconstruction error to decide 4-bit allocation means the method partly predicts its own label.

We address both limitations with \textbf{MINT} (\textbf{M}emory-\textbf{I}nformed \textbf{N}-bit \textbf{T}uning), which reformulates mixed-precision quantization as a constrained optimization problem:
\begin{equation}
    \min_{\{(b_i, g_i)\}} \sum_{i=1}^{T} \pi_i \cdot \text{NRMSE}_i(b_i, g_i) \qquad
    \text{s.t.} \quad \sum_{i=1}^{T} \text{size}_i(b_i, g_i) \leq B
    \label{eq:objective}
\end{equation}
where $b_i$ and $g_i$ are the bit-width and group size for tensor $i$, $B$ is the user's memory budget, and $\pi_i$ is a soft protection prior (\S\ref{sec:priors}). The loss function is the normalized root mean squared error (NRMSE) from the rate-distortion curve at the chosen configuration. The key insight is that \emph{both bit-width and group size} are allocation variables---prior work optimizes bit-width alone, but our evidence shows group-size selection often provides larger quality improvements than bit-width changes.

Our contributions are:
\begin{itemize}
    \item A \textbf{budget-constrained formulation} (MCKP) that jointly optimizes per-tensor bit-width and group size, replacing ad-hoc thresholds with a principled allocation under an explicit memory constraint. The greedy solver is empirically near-optimal, typically within 0.1\% of the LP relaxation bound.
    \item \textbf{Budget-targeted deployment}: users specify an exact memory target and receive a near-optimal allocation. A fitted PPL prediction curve enables quality estimation \emph{before} running the pipeline, enabling hardware-specific deployment planning (e.g., iPhone 16\,GB, RTX 4090 24\,GB, Mac M4 Pro 48\,GB).
    \item An \textbf{SQNR safety veto} providing absolute quality floors independent of relative ranking, with an empirically determined 9\,dB threshold that exploits the natural SQNR gap between 2-bit (catastrophic, max 8.7\,dB) and 3-bit (usable, min 10.4\,dB).
    \item Evidence that \textbf{group-size selection is the primary quality lever}: on Qwen3-30B-A3B, 85\% of tensors receive group size 32 instead of the default 128, and this accounts for more quality improvement than bit-width changes.
    \item A recommendation to report \textbf{median perplexity and tail statistics} alongside standard corpus perplexity, with evidence that mean PPL can produce misleading quality orderings on models with outlier sequences.
    \item Evidence that MINT \textbf{consistently outperforms calibration-based GPTQ} at matched model sizes across three MoE families (Qwen3-30B-A3B, Qwen2-57B-A14B, Mixtral-8x7B), achieving $-0.95\%$ to $-4.6\%$ lower perplexity despite being entirely data-free.
    \item \textbf{Expert-grouped allocation} for MoE models: a constraint-aware grouping strategy that treats all experts within a SwitchLinear module as a single allocation unit, correctly modeling the hardware constraint that co-located experts share quantization parameters.
    \item \textbf{Scalability to 109B parameters}: evaluation on Llama-4-Scout (109B MoE, 48 layers, 16 routed experts) demonstrates hardware-targeted deployment across memory tiers with a single analysis pass.
    \item Evidence that \textbf{quality returns diminish sharply above ${\sim}$5 average bits} and that standard downstream benchmarks \textbf{saturate at 4-bit precision}, confirming perplexity as the appropriate primary metric for quantization evaluation and providing practical deployment guidance.
\end{itemize}

% ═══════════════════════════════════════════════════════════════════════════════
\section{Related Work}
\label{sec:related}
% ═══════════════════════════════════════════════════════════════════════════════

\paragraph{Calibration-based PTQ.}
GPTQ~\cite{frantar2023gptq} applies layerwise optimal brain quantization using Hessian information from calibration data. AWQ~\cite{lin2024awq} identifies salient weight channels via activation-aware scaling. SqueezeLLM~\cite{kim2024squeezellm} combines sensitivity-based non-uniform quantization with dense-and-sparse decomposition. SpQR~\cite{dettmers2024spqr} isolates outlier weights for full-precision storage. QuIP~\cite{chee2024quip} applies random orthogonal transformations before quantization. SmoothQuant~\cite{xiao2023smoothquant} migrates quantization difficulty from activations to weights. All require calibration data.

\paragraph{Data-free quantization.}
EasyQuant~\cite{tang2024easyquant} proposes data-free PTQ via weight distribution analysis. MXQ~\cite{zhang2025mxq} uses Frobenius norm of quantization error as a single sensitivity metric. HQQ~\cite{badri2024hqq} uses half-quadratic splitting for fast uniform quantization. HIGGS~\cite{badri2025higgs} applies Hadamard rotation via the linearity theorem. These methods typically apply uniform bit-widths or coarse-grained mixed precision. MINT differs by formulating allocation as budget-constrained optimization over joint (bit-width, group-size) configurations using per-tensor rate-distortion curves.

\paragraph{Sensitivity-based mixed-precision allocation.}
LLM-MQ~\cite{li2024llmmq} formulates mixed-precision as an integer linear program but requires calibration for sensitivity estimation. SliM-LLM~\cite{huang2025slimllm} achieves group-wise mixed precision but requires calibration data. CherryQ~\cite{shang2024cherryq} identifies high-impact parameters but also requires calibration. MINT is, to our knowledge, the first method to combine data-free sensitivity estimation with constrained optimization over both bit-width and group-size at per-tensor granularity.

\paragraph{MoE quantization.}
MC-MoE~\cite{wei2025mcmoe} combines expert quantization with dynamic pruning. MoEQuant~\cite{huang2025moequant} proposes expert-balanced sampling for calibration. These highlight unique MoE challenges---uneven expert activation, sparse routing---but require calibration data. MINT's data-free approach is particularly advantageous for MoE models where calibration must cover activation patterns across hundreds of experts.

% ═══════════════════════════════════════════════════════════════════════════════
\section{Method}
\label{sec:method}
% ═══════════════════════════════════════════════════════════════════════════════

MINT operates in two phases: (1) compute a per-tensor rate-distortion curve measuring reconstruction error at each candidate (bit-width, group-size) configuration, and (2) solve a budget-constrained allocation over these curves. The entire pipeline is \emph{data-free}---it requires only the pretrained weight tensors.

\subsection{Rate-Distortion Curves}

Prior data-free methods compute reconstruction error at a single operating point (typically 4-bit), creating circularity when this error is used to decide 4-bit allocation.

MINT computes a complete \textbf{rate-distortion (RD) curve} for each weight tensor $\mathbf{W}_i$: the normalized RMSE at every valid (bit-width, group-size) configuration:
\begin{equation}
    \text{NRMSE}_i(b, g) = \frac{\text{RMS}(Q(\mathbf{W}_i; b, g) - \mathbf{W}_i)}{\text{RMS}(\mathbf{W}_i)}
\end{equation}
evaluated at configurations $\mathcal{C} = \{(2,32), (3,64), (4,32), (4,64), (4,128), (8,64), (8,128)\}$, plus full-precision (16-bit, no grouping) as the zero-distortion anchor. Here $Q(\cdot; b, g)$ denotes group-wise round-to-nearest (RTN) quantization at $b$ bits with group size $g$. The multi-point curve serves directly as the loss function for the allocator: the NRMSE at the chosen configuration $(b_i, g_i)$ is the per-tensor loss in Eq.~\ref{eq:objective}. By evaluating all candidate configurations, MINT avoids the circularity of single-point methods.

\subsection{SQNR Safety Veto}

We compute the signal-to-quantization-noise ratio (SQNR) at each configuration:
\begin{equation}
    \text{SQNR}_i(b, g) = 10 \cdot \log_{10}\!\left(\frac{\|\mathbf{W}_i\|_F^2}{\|\mathbf{W}_i - Q(\mathbf{W}_i; b, g)\|_F^2}\right) \;\text{dB}
\end{equation}

Any configuration with $\text{SQNR} < \tau_{\text{floor}}$ (default $\tau_{\text{floor}} = 9$\,dB) is \textbf{excluded} from the allocator's candidate set $\mathcal{C}_i$. This provides an \emph{absolute} safety floor: regardless of budget pressure or relative ranking, tensors will never be assigned configurations that introduce catastrophic distortion. The 9\,dB threshold is empirically determined (\S\ref{sec:sqnr_validation}): it sits in the natural gap between 2-bit SQNR (max ${\sim}$8.7\,dB) and 3-bit SQNR (min ${\sim}$10.4\,dB), blocking catastrophic 2-bit allocations while permitting usable 3-bit configurations.

\subsection{Soft Protection Priors}
\label{sec:priors}

Rather than binary hard-coded protection rules (``embeddings stay 16-bit''), MINT uses \textbf{soft multiplicative priors} $\pi_i$ that enter the optimization objective (Eq.~\ref{eq:objective}). A prior of $\pi_i = 10$ means the allocator needs $10\times$ the quality improvement per byte to justify compressing tensor $i$ versus a default tensor.

\begin{table}[h]
\centering
\caption{Soft protection priors. Higher values make the allocator less likely to compress the tensor.}
\label{tab:priors}
\begin{tabular}{lrl}
\toprule
Tensor Type & $\pi_i$ & Rationale \\
\midrule
Embedding & 10.0 & Lookup tables; quantization corrupts rare tokens \\
LM head & 10.0 & Final projection; directly affects output distribution \\
LayerNorm / RMSNorm & $\infty$ & Tiny ($<$0.01\% of params); critical for training stability \\
MoE router & 8.0 & Controls expert routing; errors cascade \\
Vision encoder & 8.0 & Cross-modal alignment is sensitive \\
First layer & 3.0 & No error correction from prior layers \\
Last layer & 2.0 & Directly precedes output \\
Default & 1.0 & No bias \\
\bottomrule
\end{tabular}
\end{table}

Tensors with $\pi_i = \infty$ are hard-protected at 16-bit and excluded from the allocator. All other priors participate in the optimization---the allocator \emph{can} compress high-prior tensors if the budget demands it and the quality tradeoff is favorable.

\subsection{Budget-Constrained Allocation}

The allocation problem is a \textbf{Multiple-Choice Knapsack Problem (MCKP)}. Each tensor $i$ has a candidate set $\mathcal{C}_i$ (configurations surviving the SQNR veto). For each candidate $(b, g) \in \mathcal{C}_i$, the \emph{cost} is $\text{size}_i(b, g)$ and the \emph{loss} is $\pi_i \cdot \text{NRMSE}_i(b, g)$.

For quantized tensors ($b < 16$), the memory cost is:
\begin{equation}
    \text{size}_i(b, g) = \left\lceil \frac{n_i \cdot b}{8} \right\rceil + \left\lceil \frac{n_i}{g} \right\rceil \cdot 4 \;\text{bytes}
    \label{eq:size}
\end{equation}
where $n_i$ is the parameter count and the second term accounts for 16-bit scale and zero-point per group. Full-precision tensors ($b = 16$) have $\text{size}_i = 2 n_i$ bytes with no grouping overhead.

The budget $B$ is specified by the user as an absolute size (\texttt{--budget 6GB}), an average bits-per-parameter target (\texttt{--avg-bits 4.0}), or a fraction of BF16 size (\texttt{--budget-ratio 0.25}).

We provide three solvers:
\begin{enumerate}
    \item \textbf{Greedy} (default): Initialize all tensors at their lowest-cost valid configuration. Compute \emph{upgrade options}---for each tensor and each higher configuration, compute the efficiency ratio $\eta = \Delta\text{loss} / \Delta\text{size}$. Sort upgrades by $\eta$ descending and greedily apply them until the budget is exhausted. Runtime: $O(T \cdot |\mathcal{C}| \cdot \log(T \cdot |\mathcal{C}|))$, under 10\,ms for 18{,}867 tensors. The greedy solver is a heuristic; it is not guaranteed to find the global optimum, but empirically achieves solutions within 0.1\% of the LP relaxation bound on all tested models.
    \item \textbf{LP relaxation}: Continuous relaxation solved via \texttt{scipy.optimize.linprog} with HiGHS, followed by rounding. Provides an upper bound on greedy gap.
    \item \textbf{ILP}: Exact solution via \texttt{scipy.optimize.milp}. Practical for models with fewer than 1{,}000 tensors.
\end{enumerate}

\subsection{Joint Bit-Width and Group-Size Optimization}

A critical design decision in MINT is that \emph{group size is an allocation variable, not a hyperparameter}. Prior mixed-precision methods optimize bit-width while keeping group size fixed (typically 128). However, empirical evidence shows that reducing group size from 128 to 64 often yields larger perplexity improvements than bit-width changes~\cite{li2025mixllm}.

MINT's configuration space includes multiple group sizes at each bit-width: $\{(4,32), (4,64), (4,128), (8,64), (8,128)\}$. The allocator considers both the quality impact (NRMSE) and the memory cost (Eq.~\ref{eq:size}) of each option, automatically selecting finer granularity for sensitive tensors and coarser granularity for insensitive ones.

\subsection{Pipeline Summary}

\begin{algorithm}[t]
\caption{MINT Analysis and Allocation Pipeline}
\label{alg:pipeline}
\begin{algorithmic}[1]
\REQUIRE Model directory, memory budget $B$, SQNR floor $\tau$
\ENSURE Per-tensor manifest: $\{(b_i, g_i)\}$ for all tensors
\STATE \textbf{Phase 1: Rate-distortion analysis}
\FOR{each shard file in model}
    \FOR{each 2D weight tensor $\mathbf{W}_i$ with $n_i \geq 1024$}
        \STATE Compute RD curve: $\text{NRMSE}_i(b,g)$ for all $(b,g) \in \mathcal{C}$
        \STATE Compute SQNR: $\text{SQNR}_i(b,g)$ for all $(b,g) \in \mathcal{C}$
    \ENDFOR
    \STATE Free shard memory
\ENDFOR
\STATE \textbf{Phase 2: Allocate}
\STATE Compute protection priors $\pi_i$ for each tensor
\STATE Filter candidate sets: $\mathcal{C}_i \leftarrow \{(b,g) \in \mathcal{C} : \text{SQNR}_i(b,g) \geq \tau\}$
\STATE Subtract protected tensor sizes from budget $B$
\STATE Run greedy MCKP solver $\rightarrow$ allocation $\{(b_i, g_i)\}$
\STATE Write manifest JSON
\end{algorithmic}
\end{algorithm}

The pipeline (Algorithm~\ref{alg:pipeline}) processes safetensor shards sequentially, keeping peak memory well below the model size. The output manifest maps each tensor to its assigned $(b_i, g_i)$ configuration and is used as a quantization predicate during model conversion via MLX~\cite{mlx2023}.

\subsection{Expert Handling for MoE Models}

For MoE models with expert tensors stored as 3D arrays $[E, d_{\text{in}}, d_{\text{out}}]$, we analyze each expert slice separately ($E \leq 32$) or cluster experts and sample representatives ($E > 32$), using worst-case SQNR across experts for safety.

\subsubsection{Expert-Grouped Allocation}

Inference frameworks (MLX, vLLM) implement MoE expert layers as \texttt{SwitchLinear} modules where all $E$ experts within a single (layer, projection) group share the same quantization parameters. The allocator must respect this constraint: it is not possible to quantize expert 0 at 3-bit and expert 1 at 4-bit within the same module.

MINT groups expert tensors by (layer, projection) before running the MCKP solver. For each group $G = \{W^{(1)}, \ldots, W^{(E)}\}$:
\begin{itemize}
    \item \textbf{NRMSE}: parameter-weighted mean across experts: $\text{NRMSE}_G(b,g) = \sum_e \frac{n_e}{n_G} \cdot \text{NRMSE}^{(e)}(b,g)$, where $n_e$ is the parameter count of expert $e$ and $n_G = \sum_e n_e$. This is consistent with the additive global objective (Eq.~\ref{eq:objective}).
    \item \textbf{SQNR}: worst-case (minimum) across experts: $\text{SQNR}_G(b,g) = \min_e \text{SQNR}^{(e)}(b,g)$. Safety should be conservative.
    \item \textbf{Size}: sum across all experts: $\text{size}_G(b,g) = \sum_e \text{size}^{(e)}(b,g)$
\end{itemize}
This reduces the number of allocation variables from $T \cdot E$ to $T$ (e.g., 5{,}376 expert tensors $\rightarrow$ 84 groups on Qwen2-57B-A14B), while ensuring that the solver's output is directly implementable by the runtime. The parameter-weighted mean NRMSE ensures that the upgrade efficiency of expert groups scales proportionally with both quality impact and size, consistent with the additive structure of the global objective.

% ═══════════════════════════════════════════════════════════════════════════════
\section{Experiments}
\label{sec:experiments}
% ═══════════════════════════════════════════════════════════════════════════════

We evaluate MINT on six model families using WikiText-2 perplexity (test split, sequence length 2048, seed 42). All experiments run on an Apple Mac Studio with M2 Ultra (192\,GB unified memory). Analysis uses PyTorch CPU; conversion and inference use MLX~\cite{mlx2023}.

\textbf{Models:} Qwen3-8B (dense, 8.2B parameters), Qwen3-30B-A3B (MoE, 30.5B parameters, 128 routed experts), Qwen2-57B-A14B (MoE, 57.4B parameters, 64 experts/layer), Mixtral-8x7B-Instruct (MoE, 46.7B parameters, 8 experts/layer), GLM-4.7-Flash (dense, 31.2B parameters), and Llama-4-Scout-17B-16E-Instruct (MoE, 109B parameters, 16 routed + 1 shared expert/layer).

\textbf{Baselines:} BF16 (unquantized), uniform 4-bit RTN with group size 64 or 128, GPTQ-Int4 where available, and the threshold-based v1 heuristic.

\subsection{Main Results: Matched-Budget Perplexity}
\label{sec:main_results}

All comparisons are at matched or near-matched model sizes, addressing the critique that mixed-precision methods win only by spending more bits.

\begin{table}[h]
\centering
\caption{Perplexity comparison at matched and varied model sizes. All evaluations use WikiText-2 (test, seq\_len=2048, 128 sequences, seed=42). Both mean and median PPL are reported; see \S\ref{sec:mean_vs_median} for discussion of when they diverge. $\Delta$ columns are relative to BF16. $\dagger$Values from~\cite{xu2025qwen3quant}. $\ddagger$BF16 model (203\,GB) exceeds available memory; $\Delta$ computed vs uniform 4-bit.}
\label{tab:main_results}
\begin{tabular}{llrrrr}
\toprule
Model & Condition & Size (GB) & Mean PPL & Median PPL & $\Delta$Med \\
\midrule
Qwen3-8B & BF16 baseline & 15.26 & 9.727 & --- & --- \\
 & AWQ 4-bit$^\dagger$ & 4.05 & 10.50 & --- & +8.1\% \\
 & GPTQ 4-bit$^\dagger$ & 4.05 & 10.30 & --- & +6.1\% \\
 & Uniform 4-bit & 4.05 & 10.249 & --- & +5.4\% \\
 & v1 (best) & 6.05 & 10.097 & --- & +3.8\% \\
 & \textbf{MINT} & \textbf{6.00} & \textbf{10.039} & --- & \textbf{+3.2\%} \\
\midrule
Qwen3-30B-A3B & BF16 baseline & 56.87 & 8.733 & 8.789 & --- \\
(MoE) & Uniform 4-bit & 15.11 & 9.629 & --- & +10.3\% \\
 & v1 (best) & 16.73 & 8.924 & 8.974 & +2.8\% \\
 & \textbf{MINT (16\,GB)} & \textbf{16.29} & \textbf{8.930} & \textbf{8.971} & \textbf{+2.3\%} \\
 & \textbf{MINT (17\,GB)} & \textbf{17.39} & \textbf{8.858} & \textbf{8.912} & \textbf{+1.5\%} \\
 & \textbf{MINT (19\,GB)} & \textbf{19.01} & \textbf{8.782} & \textbf{8.798} & \textbf{+0.6\%} \\
\midrule
GLM-4.7-Flash & BF16 baseline & 58.16 & 11.344 & 8.706 & --- \\
 & Uniform 4-bit & 14.82 & --- & ${\sim}$11.46 & +31.6\% \\
 & v1 (best) & 15.92 & 9.930 & 9.084 & +4.3\% \\
 & \textbf{MINT} & \textbf{15.82} & \textbf{9.427} & \textbf{9.210} & \textbf{+5.8\%} \\
\midrule
Llama-4-Scout$^\ddagger$ & BF16 baseline & ${\sim}$203 & \multicolumn{2}{c}{\emph{exceeds 192\,GB}} & --- \\
(109B MoE) & MINT (no safety) & 34.62 & 23.577 & 23.714 & +198\% \\
 & \textbf{MINT (min-safe)} & \textbf{46.93} & \textbf{8.675} & \textbf{8.786} & \textbf{+9.8\%} \\
 & \textbf{MINT (50\,GB)} & \textbf{51.98} & \textbf{7.980} & \textbf{8.284} & \textbf{+1.0\%} \\
 & Uniform 4-bit & 56.9 & 7.899 & --- & --- \\
 & v1 (best) & 59.5 & 7.628 & --- & $-$3.4\% \\
 & \textbf{MINT (64\,GB)} & \textbf{58.03} & \textbf{7.703} & \textbf{8.070} & \textbf{$-$2.5\%} \\
 & \textbf{MINT (192\,GB)} & \textbf{163.24} & \textbf{7.359} & \textbf{7.691} & \textbf{$-$6.8\%} \\
\bottomrule
\end{tabular}
\end{table}

\paragraph{Qwen3-8B (dense).} MINT achieves PPL 10.039 at 6.00\,GB, outperforming the best v1 configuration (PPL 10.097 at 6.05\,GB) by 0.6\%. The improvement is attributable to joint group-size optimization: MINT assigns group size 32 to 29 tensors, group size 64 to 46, and group size 128 to 101, directing finer granularity to the most sensitive tensors. MINT closes 41\% of the gap between uniform 4-bit and BF16.

\paragraph{Comparison with calibration-based methods (dense).} Published results from an independent evaluation~\cite{xu2025qwen3quant} report GPTQ (PPL 10.30) and AWQ (PPL 10.50) for Qwen3-8B at uniform 4-bit (g128). MINT (PPL 10.039) outperforms both \emph{despite requiring no calibration data}: $-2.5\%$ versus GPTQ and $-4.4\%$ versus AWQ (Table~\ref{tab:main_results}). MINT's advantage comes at a size cost (6.00\,GB vs 4.05\,GB), but this is precisely the budget-targeted tradeoff MINT is designed for: given more memory, it delivers better quality than calibration-based methods achieve with less.

\paragraph{Qwen3-30B-A3B (MoE).} This model demonstrates MINT's budget-targeting capability. At three different user-specified budgets, MINT achieves 100\% budget utilization with monotonically improving quality:
\begin{itemize}
    \item At 16.29\,GB: PPL 8.930, matching v1's quality (8.924) at \textbf{2.6\% less storage}
    \item At 17.39\,GB: PPL 8.858, surpassing v1 by 0.7\%
    \item At 19.01\,GB: PPL 8.782, within \textbf{0.6\% of BF16} at just 33\% of full size
\end{itemize}
The greedy solver allocated 18{,}867 tensors in under 20\,ms. The SQNR safety veto blocked all 2-bit configurations on this model (max SQNR: 8.1\,dB, below the 9\,dB floor), catching unsafe allocations that v1's heuristic thresholds allowed (4\% of parameters at 2-bit with SQNR below 8\,dB). 3-bit configurations (min SQNR: 9.5\,dB) pass the veto and are available to the allocator at tighter budgets.

\paragraph{GLM-4.7-Flash.} MINT achieves median PPL 9.210 at 15.82\,GB, closing 82\% of the gap between uniform 4-bit and BF16. On this model, mean PPL is misleading: BF16 exhibits 5 catastrophic outlier sequences (PPL 25{,}000--81{,}000), inflating the mean to 11.344, while quantized models stabilize these sequences (see \S\ref{sec:mean_vs_median}).

\paragraph{Llama-4-Scout (109B MoE).} This model---the largest in our evaluation at 109 billion parameters with 48 MoE layers---demonstrates MINT's budget-targeting capability at extreme scale. The BF16 model (${\sim}$203\,GB) exceeds our 192\,GB test system, making data-free quantization essential rather than optional. MINT produces two hardware-targeted variants from a single analysis pass:
\begin{itemize}
    \item \textbf{64\,GB target}: 58.03\,GB, PPL 7.703, fits a 64\,GB Mac Studio with 6\,GB headroom for KV cache and OS. Outperforms uniform 4-bit (PPL 7.899) by 2.5\%.
    \item \textbf{192\,GB target}: 163.24\,GB, PPL 7.359, fits a 192\,GB Mac Studio with 29\,GB headroom. The allocator keeps most tensors at BF16 and selectively applies 8-bit to the most compressible expert projections.
\end{itemize}
The SQNR safety veto blocks all 2-bit configurations on this architecture (SQNR $<$ 8.7\,dB), while 3-bit configurations (SQNR 10.4--14.6\,dB) pass the 9\,dB floor. This enables a \textbf{minimum safe model of 46.9\,GB} (all 3-bit, PPL 8.675)---17\% smaller than the all-4-bit floor with acceptable quality. A 2-bit model (34.6\,GB) was also tested and confirmed catastrophic (PPL 23.6), validating the veto.

\subsection{Joint Bit-Width and Group-Size Allocation}
\label{sec:allocation_detail}

Table~\ref{tab:allocation} shows MINT's allocation decisions at the 19\,GB budget on Qwen3-30B-A3B, revealing how the allocator exploits group size as a quality lever.

\begin{table}[h]
\centering
\caption{MINT per-tensor allocation at 19\,GB budget on Qwen3-30B-A3B. The allocator overwhelmingly selects group size 32 for 4-bit tensors (85\% of all tensors), trading scale/bias overhead for quantization accuracy.}
\label{tab:allocation}
\begin{tabular}{lrrl}
\toprule
Config $(b, g)$ & Tensors & \% & Role \\
\midrule
(4, 32) & 15{,}908 & 85.2\% & Default: small groups maximize accuracy \\
(4, 64) & 9 & $<$0.1\% & Transitional \\
(4, 128) & 96 & 0.5\% & Low-sensitivity tensors \\
(8, 128) & 2{,}612 & 14.0\% & Sensitive: first/last layer attention, MoE experts \\
(8, 64) & 1 & $<$0.1\% & \\
(16, 0) & 241 & 1.3\% & Hard-protected: norms, biases \\
\midrule
\textbf{Total} & 18{,}867 & & \\
\bottomrule
\end{tabular}
\end{table}

The dominant pattern is striking: \textbf{85\% of tensors receive 4-bit quantization with group size 32}, not the conventional group size 128. The allocator discovers that reducing group size from 128 to 32 trades 0.125 bytes/parameter of scale/bias overhead for substantially better per-group quantization resolution. This overhead is funded by eliminating all 16-bit allocations for quantizable tensors---v1 allocated 5.6\% of parameters to 16-bit, but the knapsack determines this is wasteful: those bytes produce more quality when redistributed as finer group sizes across 4-bit tensors.

The 8-bit allocations (14\%) concentrate on first-layer attention tensors (prior $\pi = 3.0$), last-layer attention (prior $\pi = 2.0$), and the most sensitive MoE expert weights. Embeddings and the LM head remain at 4-bit despite their prior of $\pi = 10.0$---they are too large for 8-bit upgrades to be efficient.

\subsection{SQNR Safety Veto Validation}
\label{sec:sqnr_validation}

To determine the optimal SQNR floor, we conducted a systematic sweep on Llama-4-Scout, measuring perplexity at different floor values. The SQNR distribution reveals a natural gap between bit levels that guides floor selection.

\subsubsection{SQNR Distribution Analysis}

Across both Llama-4-Scout (691 2D tensors) and Qwen3-30B-A3B (18{,}674 2D tensors), SQNR values cluster into well-separated bands by bit-width:

\begin{table}[h]
\centering
\caption{SQNR distribution (dB) across all 2D tensors on Llama-4-Scout. A natural gap of ${\sim}$2\,dB separates 2-bit (max 8.7\,dB) from 3-bit (min 10.4\,dB). Setting the floor at 9\,dB exploits this gap: all 2-bit configurations are vetoed while all 3-bit configurations are permitted.}
\label{tab:sqnr_dist}
\begin{tabular}{lrrrrrrr}
\toprule
Config & Min & P5 & Median & P95 & Max & $<$9\,dB & $<$15\,dB \\
\midrule
(2, 32) & 5.1 & 7.2 & 8.0 & 8.1 & 8.7 & 691 & 691 \\
(2, 64) & 2.5 & 5.6 & 6.8 & 6.9 & 7.2 & 691 & 691 \\
(3, 64) & 10.4 & 13.0 & 14.2 & 14.3 & 14.6 & 0 & 691 \\
(4, 32) & 19.4 & 21.3 & 22.0 & 22.1 & 22.8 & 0 & 0 \\
(4, 64) & 17.0 & 19.6 & 20.8 & 20.9 & 21.3 & 0 & 0 \\
(4, 128) & 15.1 & 18.2 & 19.8 & 20.0 & 20.0 & 0 & 0 \\
(8, 64) & 41.6 & 44.3 & 45.4 & 45.6 & 45.9 & 0 & 0 \\
\bottomrule
\end{tabular}
\end{table}

The same pattern holds on Qwen3-30B-A3B: 2-bit max SQNR is 8.1\,dB, 3-bit min SQNR is 9.5\,dB. The gap between 2-bit and 3-bit is consistent across architectures.

\subsubsection{Floor Sweep: Where Quality Breaks}

We evaluated Llama-4-Scout at four SQNR floor settings to determine where model quality transitions from acceptable to catastrophic:

\begin{table}[h]
\centering
\caption{SQNR floor sweep on Llama-4-Scout. The quality cliff is between 2-bit and 3-bit (PPL triples), not between 3-bit and 4-bit (+12.5\%). The 9\,dB floor exploits the natural SQNR gap to permit 3-bit while blocking catastrophic 2-bit.}
\label{tab:sqnr}
\begin{tabular}{lrrrrl}
\toprule
SQNR Floor & Avg Bits & Size (GB) & PPL (std) & PPL (med) & Effect \\
\midrule
0 dB (no safety) & 2.00 & 34.62 & 23.577 & 23.714 & Catastrophic (+198\%) \\
\textbf{9 dB (default)} & \textbf{3.00} & \textbf{46.93} & \textbf{8.675} & \textbf{8.786} & \textbf{Usable (+9.8\%)} \\
9 dB + 50\,GB budget & 3.48 & 51.98 & 7.980 & 8.284 & Good (+1.0\%) \\
15 dB (conservative) & 4.00 & 56.16 & 7.709 & 8.076 & Best ($-$2.4\%) \\
15 dB + 58\,GB budget & 4.01 & 58.03 & 7.703 & 8.070 & Best ($-$2.5\%) \\
\bottomrule
\end{tabular}
\end{table}

The results reveal two key findings:

\paragraph{The real cliff is between 2-bit and 3-bit.} PPL jumps from 8.675 (3-bit) to 23.577 (2-bit)---a 2.7$\times$ degradation. In contrast, 3-bit to 4-bit is only a 12.5\% increase (8.675 $\rightarrow$ 7.709). The 15\,dB floor used in earlier experiments was overly conservative: it blocked all 3-bit configurations despite them producing acceptable quality.

\paragraph{The 9\,dB floor exploits a natural gap.} On both Llama-4-Scout and Qwen3-30B-A3B, 2-bit SQNR peaks at ${\sim}$8.7\,dB while 3-bit SQNR starts at ${\sim}$10.4\,dB. A 9\,dB floor sits cleanly in this gap, vetoing all 2-bit configurations (proven catastrophic) while permitting all 3-bit configurations (proven usable). This reduces the minimum safe model size by \textbf{9.2\,GB} (53.75 $\rightarrow$ 44.27\,GB estimated, 56.16 $\rightarrow$ 46.93\,GB on disk) with only a 12.5\% PPL penalty versus the conservative floor.

\paragraph{Mixed 3/4/8-bit fills the gap.} At the 50\,GB budget with 9\,dB floor, the allocator produces a mixed allocation (88\% 3-bit, 2.6\% 4-bit, 9.1\% 8-bit) that achieves PPL 7.980---only 1.0\% worse than uniform 4-bit at 4\,GB smaller. This intermediate point was inaccessible with the 15\,dB floor.

\subsection{Budget-Targeted Deployment}
\label{sec:budget_targeting}

A unique capability of MINT's constrained formulation is \textbf{hardware-targeted quantization}: the user specifies an exact memory budget matching their deployment target, and MINT produces a near-optimal allocation for that constraint. Table~\ref{tab:budget_curve} shows the full quality--size tradeoff curve for Qwen3-30B-A3B, and Table~\ref{tab:scout_budget} demonstrates the same capability on Llama-4-Scout at extreme scale.

\begin{table}[h]
\centering
\caption{PPL vs.\ budget for Qwen3-30B-A3B. The curve enables deployment planning: given a memory target, predict quality before running MINT.}
\label{tab:budget_curve}
\begin{tabular}{rrrrl}
\toprule
Budget & Model Size & Mean PPL & Median PPL & Deployment Target \\
\midrule
15.1 GB & 15.11 GB & 9.629 & --- & \emph{(uniform 4-bit floor)} \\
15.3 GB & 16.13 GB & 8.970 & 9.020 & iPhone 16 Pro (16\,GB) \\
15.5 GB & 16.29 GB & 8.930 & 8.971 & \\
16.7 GB & 17.39 GB & 8.858 & 8.912 & \\
19.2 GB & 19.01 GB & 8.782 & 8.798 & \\
20.0 GB & 19.32 GB & 8.784 & 8.803 & RTX 4070 (20\,GB headroom) \\
25.0 GB & 27.39 GB & 8.760 & 8.779 & RTX 4090 (24\,GB) \\
30.0 GB & 30.75 GB & 8.657 & 8.684 & Mac M4 Pro (36\,GB) \\
--- & 56.87 GB & 8.733 & 8.789 & \emph{BF16 (unquantized)} \\
\bottomrule
\end{tabular}
\end{table}

Several observations emerge from this curve:

\paragraph{Steep initial gains.} The first 1\,GB above the 4-bit floor (15.1 $\rightarrow$ 16.1\,GB) produces a \textbf{6.8\% PPL improvement} (9.629 $\rightarrow$ 8.970). This is because the additional budget funds selective 8-bit upgrades for the most sensitive tensors and finer group sizes for 4-bit tensors. Threshold-based methods cannot exploit this---they produce a fixed allocation regardless of budget.

\paragraph{Diminishing returns.} Beyond ${\sim}$19\,GB, additional budget yields marginal returns: the 19 $\rightarrow$ 30\,GB increase improves mean PPL by only 1.4\%. The 30\,GB model (mean PPL 8.657) appears to outperform BF16 (mean PPL 8.733), but controlled experiments (\S\ref{sec:quant_noise}) show this is a distributional artifact rather than genuine regularization.

\paragraph{Fitted prediction curve.} We fit an inverse-power model to the measured points as a local empirical interpolation:
\begin{equation}
    \text{PPL}(B) = 8.371 + \frac{0.494}{(B - 15.099)^{0.135}}
    \label{eq:ppl_curve}
\end{equation}
with RMSE = 0.025 PPL over the measured range $B \in [15.2, 30.0]$\,GB (95\% CI: $\pm$0.046 PPL). This enables practitioners to estimate output quality for budget targets within this range \emph{before} running the pipeline. Key thresholds: within +1\% of BF16 at 17.3\,GB, within +2\% at 15.7\,GB. We caution that the parametric form is not physically meaningful outside the measured range---in particular, the fitted asymptote lies below BF16 PPL, which is an artifact of the functional form rather than a meaningful prediction.

\paragraph{Replication on Mixtral-8x7B.} Table~\ref{tab:mixtral_budget} shows the budget curve for Mixtral-8x7B, confirming the same pattern on a different MoE architecture (8 experts per layer vs Qwen3's 128).

\begin{table}[h]
\centering
\caption{PPL vs.\ budget for Mixtral-8x7B-Instruct. The same diminishing-returns pattern and steep initial gains observed on Qwen3-30B-A3B replicate on a different MoE architecture. BF16 baseline (${\sim}$87\,GB) exceeds practical deployment targets.}
\label{tab:mixtral_budget}
\begin{tabular}{rrrrl}
\toprule
Budget & Model Size & Mean PPL & Median PPL & Note \\
\midrule
--- & 24.5 GB & 4.387 & 4.413 & \emph{(uniform 4-bit baseline)} \\
min-safe & 19.4 GB & 4.926 & 4.993 & All 3-bit (${\sim}$3.1 avg bits) \\
24.5 GB & 24.5 GB & 4.264 & 4.266 & \textbf{Matched-size: $-$2.8\% vs uniform} \\
32.0 GB & 31.9 GB & 4.218 & 4.259 & ${\sim}$5.0 avg bits \\
39.0 GB & 39.2 GB & 4.198 & 4.241 & ${\sim}$6.7 avg bits \\
--- & 46.2 GB & 4.174 & 4.216 & \emph{(uniform 8-bit g64 reference)} \\
\bottomrule
\end{tabular}
\end{table}

The Mixtral curve mirrors Qwen3-30B: steep gains in the first GB above the 4-bit floor (4.387 $\rightarrow$ 4.264, $-2.8\%$ at matched size), followed by rapid flattening. At 31.9\,GB (${\sim}$5 avg bits), MINT achieves PPL 4.218---within 1.1\% of the 8-bit reference (4.174 at 46.2\,GB) at 31\% less storage. The pattern of diminishing returns above ${\sim}$5 bits is consistent across both architectures.

\paragraph{Scaling to 109B parameters.} Table~\ref{tab:scout_budget} demonstrates budget targeting on Llama-4-Scout, a model too large to run in BF16 on any single consumer device. From a single rate-distortion analysis pass, MINT generates allocations for different hardware tiers.

\begin{table}[h]
\centering
\caption{MINT budget-targeted deployment of Llama-4-Scout (109B MoE). The same analysis produces allocations for different hardware targets. BF16 (${\sim}$203\,GB) exceeds all single-device targets.}
\label{tab:scout_budget}
\begin{tabular}{rrrrl}
\toprule
Budget & Model Size & Mean PPL & Median PPL & Deployment Target \\
\midrule
--- & 34.62 GB & 23.577 & 23.714 & \emph{(no safety --- catastrophic)} \\
min-safe & 46.93 GB & 8.675 & 8.786 & \emph{min safe (all 3-bit)} \\
50 GB & 51.98 GB & 7.980 & 8.284 & Mac Studio M4 (48\,GB headroom) \\
56.9 GB & 56.9 GB & 7.899 & --- & \emph{(uniform 4-bit baseline)} \\
64 GB & 58.03 GB & 7.703 & 8.070 & Mac Studio M4 Max (64\,GB) \\
192 GB & 163.24 GB & 7.359 & 7.691 & Mac Studio M2 Ultra (192\,GB) \\
--- & ${\sim}$203 GB & \multicolumn{2}{c}{\emph{exceeds 192\,GB}} & \emph{BF16 (unquantized)} \\
\bottomrule
\end{tabular}
\end{table}

The Scout budget curve spans a 4.7$\times$ size range (35--163\,GB) and reveals three regimes. Below the 9\,dB safety floor, quality collapses (PPL 23.6 at 2-bit). Above the floor, the \texttt{--min-safe} mode produces the smallest viable model (46.9\,GB, all 3-bit, PPL 8.675). With additional budget, the allocator upgrades the most sensitive tensors: the 50\,GB mixed model (88\% 3-bit, 9\% 8-bit) achieves PPL 7.980---only 1.0\% worse than uniform 4-bit at 5\,GB smaller. The 64\,GB and 192\,GB allocations use progressively more 4-bit and 8-bit to push quality further. The 3.5$\times$ size difference between the minimum safe model and the 192\,GB model (47 vs 163\,GB) produces only a 1.9$\times$ PPL difference (8.675 vs 7.359), demonstrating diminishing returns---practitioners with constrained hardware sacrifice surprisingly little quality.

\subsection{Matched-Size Comparison with GPTQ}
\label{sec:gptq_comparison}

To rigorously evaluate MINT against calibration-based quantization, we compare against GPTQ~\cite{frantar2023gptq} at \textbf{exactly matched model sizes} across three MoE model families spanning different architectures and scales. For each family, we obtain an official GPTQ-Int4 quantized model, evaluate its perplexity, then run MINT's pipeline (rate-distortion curves $\rightarrow$ MCKP allocation $\rightarrow$ conversion) targeting the same byte budget.

A critical engineering detail for MoE models is \textbf{expert-grouped allocation}: frameworks such as MLX implement MoE experts via \texttt{SwitchLinear} modules where all experts within a layer share the same quantization parameters (bit-width and group size). MINT's allocator groups expert tensors by (layer, projection) and uses the \emph{parameter-weighted mean NRMSE} across experts within each group, ensuring that the allocation is implementable by the inference runtime and that upgrade efficiency scales consistently with the additive global objective.

\begin{table}[h]
\centering
\caption{Matched-size comparison: MINT vs GPTQ across three MoE families. All evaluations use WikiText-2 (test, seq\_len=2048, seed=42). MINT outperforms GPTQ on all three families despite being entirely data-free. $^\dagger$Mixtral GPTQ uses per-channel quantization (group\_size=$-1$); original compressed size is 22\,GB but must be dequantized to fp16 (87\,GB) for evaluation since MLX does not support per-channel quantization. MINT comparison is at the uniform 4-bit size (24.5\,GB).}
\label{tab:gptq_comparison}
\begin{tabular}{llrrrr}
\toprule
Model & Method & Size (GB) & Mean PPL & Med PPL & $\Delta$ vs GPTQ \\
\midrule
Qwen3-30B-A3B & GPTQ Int4 (g128) & 16.0 & 9.122 & 9.160 & --- \\
(MoE, 30B) & \textbf{MINT} & \textbf{16.1} & \textbf{8.970} & \textbf{9.020} & \textbf{$-$1.5\%} \\
\midrule
Qwen2-57B-A14B & GPTQ Int4 (g128) & 29.9 & 6.390 & 6.396 & --- \\
(MoE, 57B) & \textbf{MINT} & \textbf{29.9} & \textbf{6.329} & \textbf{6.356} & \textbf{$-$0.6\%} \\
\midrule
Mixtral-8x7B$^\dagger$ & GPTQ per-ch & 87.0$^\dagger$ & 4.608 & 4.640 & --- \\
(MoE, 47B) & Uniform 4-bit (g64) & 24.5 & 4.471 & 4.461 & --- \\
 & \textbf{MINT} & \textbf{24.5} & \textbf{4.264} & \textbf{4.266} & \textbf{$-$4.6\%} \\
\bottomrule
\end{tabular}
\end{table}

\paragraph{Consistent wins across architectures.} MINT outperforms GPTQ on all three model families, with improvements ranging from $-0.6\%$ to $-4.6\%$ median PPL at matched sizes. These are different MoE architectures (Qwen2-MoE with 64 experts per layer, Qwen3 with shared expert + 128 routed experts, Mixtral with 8 experts) ranging from 30B to 57B total parameters. The consistency across architectures is notable because GPTQ uses Hessian-based calibration with activation data while MINT uses only the weight tensors.

\paragraph{Expert grouping is essential.} Without expert grouping, the allocator treats each expert tensor independently and assigns different bit-widths to experts within the same SwitchLinear module. Since the inference runtime requires uniform quantization within each module, the bridge must aggregate these decisions (typically via mode), creating a mismatch between planned and actual allocation. On Qwen2-57B-A14B, this mismatch inflated the model from 29.9\,GB (planned) to 30.2\,GB (actual) and degraded PPL from 6.329 to 7.147. Expert grouping with parameter-weighted mean NRMSE eliminates this mismatch entirely: the allocator's 29.9\,GB estimate matches the actual output size to within 0.01\,GB.

\paragraph{Why does data-free beat calibration?} We hypothesize three factors: (1) \textbf{joint group-size optimization}: MINT explores a configuration space of 8 (bits, group\_size) combinations per tensor, while GPTQ uses a fixed group size; (2) \textbf{budget-constrained allocation}: MINT's MCKP solver targets a specific memory budget, while GPTQ applies uniform bit-width without budget awareness; and (3) \textbf{MoE calibration difficulty}: calibrating quantization across hundreds of experts requires activating each expert sufficiently, which narrow calibration sets may fail to do. MINT's data-free approach avoids this coverage problem entirely.

\subsection{Mean vs.\ Median Perplexity}
\label{sec:mean_vs_median}

Standard perplexity evaluation computes the geometric mean loss across all evaluation sequences. We observe that this metric can be misleading:

\begin{table}[h]
\centering
\caption{Mean vs.\ median perplexity reveals different quality orderings. GLM-4.7-Flash is the extreme case, but the effect appears across models.}
\label{tab:mean_median}
\begin{tabular}{llrrr}
\toprule
Model & Condition & Mean PPL & Median PPL & Outliers ($>$100) \\
\midrule
Qwen3-30B & MINT (19\,GB) & 8.782 & 8.798 & 0 \\
 & MINT (16\,GB) & 8.930 & 8.971 & 0 \\
 & v1 (best) & 8.924 & 8.974 & 0 \\
\midrule
GLM-4.7 & BF16 & 11.344 & 8.706 & 5 \\
 & v1 & 9.930 & 9.084 & 4 \\
 & \textbf{MINT} & \textbf{9.427} & \textbf{9.210} & \textbf{0} \\
\midrule
Llama-4-Scout & Uniform 4-bit & 7.899 & --- & --- \\
 & \textbf{MINT (64\,GB)} & \textbf{7.703} & \textbf{8.070} & \textbf{0} \\
 & \textbf{MINT (192\,GB)} & \textbf{7.359} & \textbf{7.691} & \textbf{0} \\
\bottomrule
\end{tabular}
\end{table}

On GLM-4.7-Flash, mean PPL gives a completely \textbf{inverted quality ordering}: BF16 appears worst (11.344) and MINT best (9.427). Median PPL gives the correct ranking: BF16 best (8.706), v1 second (9.084), MINT third (9.210). The inversion is caused by 5 catastrophic outlier sequences where BF16 produces PPL values of 25{,}000--81{,}000, while quantization noise acts as implicit regularization that stabilizes these pathological sequences to PPL 100--360.

On Qwen3-30B-A3B, the mean and median are consistent in their ordering but differ in magnitude: MINT (16\,GB) has mean 8.930 but median 8.971. The 0.04 difference matters when comparing against v1's mean of 8.924---mean PPL suggests MINT is worse, while median PPL (8.971 vs 8.974) shows MINT is marginally better at smaller size.

\textbf{Recommendation:} Future quantization evaluations should report standard corpus perplexity (token-weighted cross-entropy) as the primary metric, supplemented by median per-sequence PPL, tail percentiles (P95, P99), and outlier counts as robustness diagnostics. When the per-sequence loss distribution is heavy-tailed, mean per-sequence PPL can produce misleading orderings; reporting both gives a more complete picture of quantization quality.

\subsection{Analysis Efficiency}

\begin{table}[h]
\centering
\caption{MINT pipeline timing on M2 Ultra 192\,GB.}
\label{tab:timing}
\begin{tabular}{lrrrr}
\toprule
Model & Tensors & Pass 1 & Allocation & Total \\
\midrule
Qwen3-8B & 399 & 3 min & $<$1s & ${\sim}$10 min \\
Qwen3-30B-A3B & 18{,}867 & 50 min & $<$1s & ${\sim}$54 min \\
GLM-4.7-Flash & 9{,}703 & 39 min & $<$1s & ${\sim}$44 min \\
Llama-4-Scout & ${\sim}$1{,}000 & 45 min & $<$1s & ${\sim}$50 min \\
\bottomrule
\end{tabular}
\end{table}

Pass 1 dominates runtime because it computes rate-distortion curves at 8 configurations per tensor. The allocation phase (greedy solver) completes in under 10\,ms regardless of model size. Total pipeline time is 2--5$\times$ slower than v1's single-pass approach but produces strictly better allocations due to joint optimization.

% ═══════════════════════════════════════════════════════════════════════════════
\section{Discussion}
\label{sec:discussion}
% ═══════════════════════════════════════════════════════════════════════════════

\subsection{Key Findings}
\label{sec:key_findings}

Our experiments across six model families yield several findings that extend beyond MINT's specific design and inform the broader quantization literature.

\begin{enumerate}

\item \textbf{Group size is the primary quality lever---more impactful than bit-width selection.}
On Qwen3-30B-A3B, the MCKP allocator chose 4-bit with group size 32 for 85\% of tensors instead of the conventional group size 128. The v1 heuristic kept all 4-bit tensors at g128 and spent its budget on 16-bit and 8-bit upgrades instead. The knapsack solver discovers that $4\times$ more quantization groups (g32 vs g128) at 0.125\,bytes/parameter overhead is more efficient than upgrading select tensors to 8-bit or 16-bit at 0.5--1.5\,bytes/parameter. Consequently, MINT \emph{never} assigns 16-bit to quantizable tensors---the bytes are always better spent on finer group sizes across thousands of other tensors. This validates the observation from prior work~\cite{li2025mixllm} that group size is a first-order compression variable and suggests that future quantization methods should treat group size as an allocation variable rather than a fixed hyperparameter.

\item \textbf{The SQNR safety veto exploits a natural gap between 2-bit and 3-bit quantization.}
Across all tested architectures, 2-bit SQNR peaks at ${\sim}$8.7\,dB while 3-bit SQNR starts at ${\sim}$10.4\,dB, leaving a ${\sim}$2\,dB gap where no valid configurations exist. Setting the floor at 9\,dB cleanly separates catastrophic quantization (PPL triples at 2-bit) from usable quantization (PPL degrades only 12.5\% at 3-bit). On Qwen3-30B-A3B, the v1 heuristic allocated 4.0\% of parameters to 2-bit with SQNR values of 4.8--8.1\,dB; the veto catches all of these. This natural gap appears consistent across dense and MoE architectures and may be a useful design principle for other quantization methods.

\item \textbf{Standard mean perplexity can produce misleading quality orderings.}
On GLM-4.7-Flash, mean PPL gives a \emph{completely inverted} ranking: BF16 appears worst (11.344) and MINT best (9.427), while median PPL gives the correct ordering (BF16 8.706, MINT 9.210). The inversion is caused by 5 catastrophic outlier sequences (PPL 25{,}000--81{,}000) in BF16. Controlled noise experiments on Qwen3-30B-A3B (Table~\ref{tab:noise_experiment}) confirm that apparent ``regularization'' from quantization is a distributional artifact: matched-scale Gaussian noise produces the same median shift ($-0.38\%$), and 68\% of individual sequences degrade. We recommend reporting median PPL and tail statistics alongside standard corpus perplexity.

\item \textbf{Data-free MINT consistently outperforms calibration-based GPTQ at matched model sizes.}
In matched-size comparisons across three MoE families, MINT outperforms GPTQ by $-0.6\%$ to $-4.6\%$ median PPL (\S\ref{sec:gptq_comparison}), and on dense Qwen3-8B, MINT outperforms both GPTQ ($-2.5\%$) and AWQ ($-4.4\%$). We attribute this to three factors: (1) joint group-size optimization accesses a richer configuration space than fixed-group GPTQ; (2) the MCKP solver targets quality maximization under an explicit budget constraint; and (3) for MoE models, narrow calibration sets may fail to cover activation patterns across hundreds of experts.

\item \textbf{Quality returns diminish sharply above ${\sim}$5 average bits.}
On Qwen3-30B-A3B, MINT at 21.2\,GB (${\sim}$5.1 avg bits) achieves PPL 8.756---within 0.1\% of the uniform 8-bit reference (PPL 8.765 at 29.3\,GB) at \textbf{28\% less storage}. On Mixtral-8x7B, the same pattern holds: MINT at 31.9\,GB (${\sim}$5.0 avg bits) is within 1.1\% of the 8-bit reference at 31\% less storage. Beyond ${\sim}$5 bits, additional bytes produce negligible quality improvement; users are better served by allocating extra memory to context length or batch size.

\item \textbf{Downstream benchmarks confirm quality preservation but are insensitive to quantization differences.}
We evaluate on ARC-Challenge (25-shot) and Winogrande (5-shot) across six budget levels per model using \texttt{lm-evaluation-harness}~\cite{eval-harness} (Table~\ref{tab:benchmarks}). MINT matches or outperforms uniform 4-bit on all downstream metrics at matched sizes (Mixtral: $+0.5/+1.3$ pp on ARC-C/WG). However, \textbf{PPL is 6--7$\times$ more sensitive} than downstream accuracy: Qwen3-30B PPL spans 10.49 to 8.73 (20\% range) from min-safe to 8-bit, while ARC-C spans only 67.2--69.9 (2.7 percentage points). Mixtral shows the same pattern: 18\% PPL range vs.\ 3.4 pp on ARC-C. Standard reasoning benchmarks effectively saturate at ${\sim}$4-bit precision and cannot discriminate between quantization quality levels above that threshold. This confirms PPL as the appropriate primary metric for quantization evaluation, while downstream benchmarks serve as a sanity check that quality is preserved.

\item \textbf{Budget-constrained formulation enables hardware-targeted deployment from a single analysis pass.}
The greedy solver achieves 100\% budget utilization across all experiments. Users specify their exact available memory---16\,GB for iPhone, 24\,GB for RTX 4090, 64--192\,GB for Mac---and MINT delivers the best achievable quality for that constraint. Combined with the PPL prediction curve (Eq.~\ref{eq:ppl_curve}), users can estimate output quality \emph{before} running the pipeline. The Llama-4-Scout results demonstrate this at extreme scale: the same 109B model (${\sim}$203\,GB in BF16, too large for any consumer device) is deployed at 58\,GB for a 64\,GB device and 163\,GB for a 192\,GB device, with only a 4.7\% quality difference despite a 2.8$\times$ size difference.

\end{enumerate}

\subsection{Supporting Analysis}

\paragraph{Does quantization act as regularization?}
\label{sec:quant_noise}
At 30.75\,GB (97.5\% 8-bit), MINT achieves mean PPL 8.657---apparently below BF16's 8.733. To test whether this reflects genuine regularization, we conducted controlled experiments on Qwen3-30B-A3B comparing BF16, uniform 8-bit quantization at two group sizes, uniform 6-bit, and Gaussian noise injection calibrated to match 8-bit quantization error magnitude ($\sigma$ matched per tensor, 385 tensors perturbed, average SNR 44.7\,dB).

\begin{table}[h]
\centering
\caption{Controlled noise experiments on Qwen3-30B-A3B. Mean and median PPL diverge in opposite directions under all noise conditions, including unstructured Gaussian noise.}
\label{tab:noise_experiment}
\begin{tabular}{lrrrr}
\toprule
Condition & Size (GB) & Mean PPL & Median PPL & $\Delta$Med \\
\midrule
BF16 (baseline) & 56.87 & 8.733 & 8.789 & --- \\
Gaussian noise ($1\times$) & 56.87 & 8.742 & 8.755 & $-$0.38\% \\
Uniform 8-bit g64 & 30.21 & 8.750 & 8.765 & $-$0.27\% \\
Uniform 8-bit g128 & 29.32 & 8.769 & 8.776 & $-$0.15\% \\
Uniform 6-bit g64 & 23.11 & 8.765 & 8.798 & +0.10\% \\
\bottomrule
\end{tabular}
\end{table}

All quantized conditions produce \emph{higher} mean PPL than BF16, ruling out genuine regularization. The median improvement ($-0.15\%$ to $-0.27\%$) is replicated by matched-scale Gaussian noise ($-0.38\%$), confirming it is a general noise phenomenon: BF16 has a left tail of sequences with unusually low PPL, and any perturbation---structured or not---partially disrupts these, pulling median down while pushing mean up. Per-sequence analysis shows 68\% of sequences degrade under 8-bit quantization, uniformly distributed across PPL quartiles, ruling out selective stabilization.

\paragraph{When the allocator disagrees with heuristics.}
On GLM-4.7-Flash, the v1 heuristic achieved slightly better median PPL (9.084 vs 9.210) by assigning 141 tensors to 8-bit. MINT's allocator determined that within the 16\,GB budget, all quantizable tensors are best at 4-bit. This reveals a limitation of NRMSE as the sole loss function: it may undervalue certain tensor-level quality improvements important for downstream perplexity. Incorporating activation-aware sensitivity (e.g., Fisher information) as a per-tensor importance weight could address this gap.

\subsection{Downstream Benchmarks}
\label{sec:downstream_benchmarks}

To validate that MINT's perplexity gains translate to downstream task performance---and to characterize how standard benchmarks compare with PPL as quality measures---we evaluate on ARC-Challenge (25-shot, normalized accuracy) and Winogrande (5-shot) using \texttt{lm-evaluation-harness}~\cite{eval-harness} via the MLX backend. We test six configurations per model spanning the full 3-bit to 8-bit range: min-safe (all 3-bit), MINT at the default budget, uniform 4-bit at matched size, MINT at $+30\%$ and $+60\%$ budgets, and uniform 8-bit as reference.

\begin{table}[h]
\centering
\caption{Perplexity and downstream accuracy across budget levels. PPL varies 18--20\% from min-safe to 8-bit reference, while ARC-C varies only 2.7--3.4 pp---confirming that downstream accuracy is preserved at 4-bit and above, and that PPL is the more discriminating metric for quantization evaluation.}
\label{tab:benchmarks}
\begin{tabular}{llrrrr}
\toprule
Model & Config & Size (GB) & PPL & ARC-C & WG \\
\midrule
Qwen3-30B & Min-safe & 13.4 & 10.494 & 67.2 & 71.3 \\
 & \textbf{MINT} & \textbf{16.3} & \textbf{8.972} & \textbf{69.5} & \textbf{69.5} \\
 & Uniform 4-bit & 16.0 & 9.098 & 69.5 & 70.3 \\
 & MINT +30\% & 21.2 & 8.756 & 69.5 & 70.1 \\
 & MINT +60\% & 26.1 & 8.771 & 69.8 & 70.5 \\
 & 8-bit ref & 29.3 & 8.765 & 69.9 & 70.9 \\
\midrule
Mixtral-8x7B & Min-safe & 19.4 & 4.926 & 67.9 & 80.2 \\
 & \textbf{MINT} & \textbf{24.5} & \textbf{4.264} & \textbf{70.5} & \textbf{81.4} \\
 & Uniform 4-bit & 24.5 & 4.387 & 70.0 & 80.1 \\
 & MINT +30\% & 31.9 & 4.218 & 71.0 & 81.7 \\
 & MINT +60\% & 39.2 & 4.198 & 71.2 & 82.1 \\
 & 8-bit ref & 46.0 & 4.174 & 71.3 & 81.9 \\
\bottomrule
\end{tabular}
\end{table}

The results show that \textbf{MINT matches or outperforms uniform 4-bit on all downstream metrics} at matched sizes (Mixtral: $+0.5/+1.3$ pp on ARC-C/WG; Qwen3: within noise). Crucially, while PPL captures a 20\% quality range from min-safe to 8-bit, ARC-C spans only 2.7 pp on Qwen3 and 3.4 pp on Mixtral. The benchmarks cannot distinguish between the MINT default, the $+30\%$ budget, the $+60\%$ budget, and the 8-bit reference---all four produce ARC-C scores within 0.4 pp on Qwen3 and 0.3 pp on Mixtral. This does \emph{not} mean the models are identical in quality---PPL clearly shows progressive improvement---but rather that \textbf{standard reasoning benchmarks saturate at ${\sim}$4-bit precision} and are too coarse to discriminate quantization quality above that level. PPL remains the appropriate primary metric for quantization evaluation; downstream benchmarks serve to confirm that quality is broadly preserved.

\subsection{Practical Guidance}
\label{sec:practical_guidance}

The scaling data yields concrete deployment recommendations:
\begin{itemize}
    \item \textbf{Default ${\sim}$4 average bits is the optimal operating point}: 97--99\% of BF16 quality at 25--30\% of the size, and MINT's per-tensor allocation provides measurable gains over uniform quantization.
    \item \textbf{Beyond ${\sim}$5 average bits offers negligible improvement}: on Qwen3-30B, 21.2\,GB (${\sim}$5 bits) matches 8-bit quality (29.3\,GB) within 0.1\% PPL---a negligible gain for 38\% more storage. Users with extra memory are better served by increased context length or batch size.
    \item \textbf{Min-safe is viable for extreme constraints}: deploying at all 3-bit costs 15--20\% PPL but produces the smallest viable model (e.g., 46.9\,GB on Llama-4-Scout vs 56.9\,GB at uniform 4-bit).
    \item \textbf{MINT's advantage is strongest at ${\sim}$4 bits}: at 8-bit all tensors quantize well and uniform allocation is nearly optimal; at 3-bit all tensors are impaired and allocation has limited room to help. The 4-bit regime is where per-tensor decisions pay off most.
\end{itemize}

\subsection{Throughput}
\label{sec:throughput}

MINT's allocator frequently selects group size 32 over the conventional 128, which increases scale/bias metadata bandwidth during inference. Table~\ref{tab:throughput} measures the impact.

\begin{table}[h]
\centering
\caption{Throughput comparison: MINT vs uniform 4-bit (g128). Apple M2 Ultra, MLX 0.30.3, seq\_len=2048, gen\_len=256, median of 3 runs. $^\dagger$Prompt penalty on MoE models is attributable to a known MLX kernel inefficiency with small group sizes (issue \#3251); we expect this gap to narrow.}
\label{tab:throughput}
\begin{tabular}{llrrrrr}
\toprule
Model & Method & Size & \multicolumn{2}{c}{Throughput (tok/s)} & $\Delta$ Gen \\
\cmidrule(lr){4-5}
 & & (GB) & Prompt & Gen & \\
\midrule
Qwen3-30B & Uniform & 17.2 & 1{,}308 & 84.4 & --- \\
 & \textbf{MINT} & \textbf{16.3} & \textbf{1{,}320} & \textbf{78.6} & $-$6.9\% \\
\midrule
Mixtral-8x7B & Uniform & 24.6 & 729 & 66.5 & --- \\
 & \textbf{MINT} & \textbf{24.5} & \textbf{332}$^\dagger$ & \textbf{57.9} & $-$12.9\% \\
\midrule
Llama-4-Scout & Uniform & 56.9 & 570 & 43.4 & --- \\
 & \textbf{MINT} & \textbf{58.0} & \textbf{299}$^\dagger$ & \textbf{37.3} & $-$14.0\% \\
\bottomrule
\end{tabular}
\end{table}

Generation throughput decreases by 7--14\% due to the additional scale/bias reads required by smaller group sizes. On Qwen3-30B, prompt processing is unaffected (MINT is slightly faster due to smaller model size), but on Mixtral and Scout the prompt path shows a 47--55\% slowdown. This traces to MLX's Metal \texttt{quantized\_matmul} kernel reading scale/bias values per-thread without threadgroup caching---a known issue we have reported (MLX issue \#3251\footnote{\url{https://github.com/ml-explore/mlx/issues/3251}}). The generation path uses a different (QMV) kernel that is less affected. We expect the prompt gap to close once the kernel is optimized; the generation overhead of 7--14\% represents the realistic steady-state cost.

\paragraph{Limitations.}
The allocator uses raw NRMSE as its sole loss function without per-tensor importance weighting. NRMSE measures weight reconstruction fidelity, but not all tensors contribute equally to downstream loss---incorporating activation-aware sensitivity (e.g., Fisher information or Hessian traces) could improve allocation quality, particularly on models like GLM-4.7-Flash where the NRMSE-based allocator disagrees with heuristic methods. We investigate several candidate sensitivity features (spectral, kurtosis, output noise amplification) in Appendix~\ref{app:features}, but in the current system these are not used by the allocator. The PPL prediction curve (Eq.~\ref{eq:ppl_curve}) is a local empirical fit valid over the measured budget range; it should not be extrapolated as a universal law. BF16 perplexity is unavailable for Llama-4-Scout as the model (${\sim}$203\,GB) exceeds our test system's 192\,GB memory. The GPTQ comparison uses publicly available GPTQ-Int4 models with default settings; comparisons against GPTQ with optimized hyperparameters (e.g., act-order with variable group sizes) could narrow the gap.

% ═══════════════════════════════════════════════════════════════════════════════
\section{Conclusion}
\label{sec:conclusion}
% ═══════════════════════════════════════════════════════════════════════════════

We presented MINT, a data-free mixed-precision quantization framework that replaces ad-hoc heuristic scoring with budget-constrained allocation over per-tensor rate-distortion curves. We evaluated MINT on six model families spanning 8B to 109B parameters, including both dense and MoE architectures. By formulating per-tensor allocation as a Multiple-Choice Knapsack Problem over joint (bit-width, group-size) configurations, MINT enables three capabilities that prior methods lack:

\begin{enumerate}
    \item \textbf{Budget-targeted deployment.} Users specify their exact memory constraint---16\,GB for iPhone, 24\,GB for RTX 4090, 64--192\,GB for Mac---and receive a near-optimal allocation. On Llama-4-Scout (109B), a single analysis pass produces deployments from 58\,GB (PPL 7.703) to 163\,GB (PPL 7.359), both outperforming uniform 4-bit. A fitted prediction curve estimates output quality before running the pipeline.
    \item \textbf{Joint group-size optimization.} Treating group size as an allocation variable (not a fixed hyperparameter) emerges as the single most impactful design decision. On Qwen3-30B-A3B, 85\% of tensors receive group size 32 instead of the conventional 128, and this produces more quality improvement than bit-width changes.
    \item \textbf{Safety guarantees with empirically validated thresholds.} The SQNR veto at 9\,dB exploits the natural gap between catastrophic 2-bit (max 8.7\,dB) and usable 3-bit (min 10.4\,dB). On Llama-4-Scout, this enables a \texttt{--min-safe} mode producing the smallest viable model (46.9\,GB, PPL 8.675) while blocking 2-bit allocations that would triple perplexity (PPL 23.6).
\end{enumerate}

On Qwen3-30B-A3B, MINT achieves +0.6\% perplexity versus BF16 at 19\,GB (33\% of full size), and matches prior heuristic quality at 2.6\% less storage. On Llama-4-Scout (109B MoE), MINT enables deployment of a model too large for BF16 on any consumer device, spanning from 46.9\,GB (min-safe, PPL 8.675) through 58\,GB (PPL 7.703) to 163\,GB (PPL 7.359)---all superior to uniform 4-bit (PPL 7.899) except the min-safe floor which trades 9.8\% quality for 17\% smaller size. Scaling experiments across budget levels reveal that quality returns diminish sharply above ${\sim}$5 average bits, and that standard downstream benchmarks (ARC-Challenge, Winogrande) saturate at 4-bit precision---varying only 2.7--3.4 percentage points across the entire 3-to-8-bit range---confirming perplexity as the appropriate primary metric for quantization evaluation. We additionally demonstrate that standard mean perplexity can produce misleading quality orderings due to outlier sequences, and recommend that quantization evaluations report median perplexity and tail statistics alongside standard corpus perplexity.

Contrary to the conventional assumption that calibration-based methods are strictly superior, MINT consistently outperforms GPTQ at matched model sizes across three MoE families (Qwen3-30B-A3B: $-1.5\%$, Qwen2-57B-A14B: $-0.6\%$, Mixtral-8x7B: $-4.6\%$ median PPL) and also outperforms both GPTQ and AWQ on dense Qwen3-8B. This consistency across six model families and two architecture types (dense and MoE) suggests that joint group-size optimization and budget-constrained allocation can provide quality gains that compensate for the absence of calibration data. The throughput cost of MINT's mixed group sizes is modest: 7--14\% generation overhead, with larger prompt-processing penalties attributable to a known framework kernel inefficiency that we expect to be resolved. For MoE models specifically, MINT's data-free approach avoids the coverage problem inherent in calibrating across hundreds of sparsely activated experts. MINT is useful both as a standalone method when calibration is impractical (very large models, proprietary data) and as a front-end for calibration-based refinement.

% ═══════════════════════════════════════════════════════════════════════════════
% References
% ═══════════════════════════════════════════════════════════════════════════════

\bibliographystyle{plain}

\begin{thebibliography}{20}

\bibitem{frantar2023gptq}
E.~Frantar, S.~Ashkboos, T.~Hoefler, and D.~Alistarh.
\newblock GPTQ: Accurate post-training quantization for generative pre-trained transformers.
\newblock In \emph{ICLR}, 2023.

\bibitem{lin2024awq}
J.~Lin, J.~Tang, H.~Tang, S.~Yang, W.-M.~Chen, W.-C.~Wang, G.~Xiao, X.~Dang, C.~Gan, and S.~Han.
\newblock AWQ: Activation-aware weight quantization for on-device LLM compression and acceleration.
\newblock In \emph{MLSys}, 2024.

\bibitem{kim2024squeezellm}
S.~Kim, C.~Hooper, A.~Gholami, Z.~Dong, X.~Li, S.~Shen, M.~W.~Mahoney, and K.~Keutzer.
\newblock SqueezeLLM: Dense-and-sparse quantization.
\newblock In \emph{ICML}, 2024.

\bibitem{dettmers2024spqr}
T.~Dettmers, R.~Svirschevski, V.~Egiazarian, D.~Kuznedelev, E.~Frantar, S.~Ashkboos, A.~Borzunov, T.~Hoefler, and D.~Alistarh.
\newblock SpQR: A sparse-quantized representation for near-lossless LLM weight compression.
\newblock In \emph{ICLR}, 2024.

\bibitem{chee2024quip}
J.~Chee, Y.~Cai, V.~Kuleshov, and C.~De~Sa.
\newblock QuIP: 2-bit quantization of large language models with guarantees.
\newblock In \emph{NeurIPS}, 2024.

\bibitem{tang2024easyquant}
J.~Tang, Y.~Liu, H.~Yin, Y.~Li, Y.~Jiang, and Q.~Tan.
\newblock EasyQuant: An efficient data-free quantization algorithm for LLMs.
\newblock \emph{arXiv preprint arXiv:2403.02775}, 2024.

\bibitem{zhang2025mxq}
F.~Zhang et al.
\newblock A mixed quantization approach for data-free quantization of LLMs.
\newblock In \emph{ICAART}, 2025.

\bibitem{badri2025higgs}
M.~Badri, F.~Tramer, and T.~Hoefler.
\newblock Pushing the limits of large language model quantization via the linearity theorem.
\newblock In \emph{NAACL}, 2025.

\bibitem{zhao2025kurtail}
Y.~Zhao et al.
\newblock KurTail: Kurtosis-based LLM quantization.
\newblock In \emph{Findings of EMNLP}, 2025.

\bibitem{li2024llmmq}
J.~Li et al.
\newblock LLM-MQ: Mixed-precision quantization for efficient LLM deployment.
\newblock Technical report, Tsinghua University, 2024.

\bibitem{li2025mixllm}
Y.~Li et al.
\newblock MixLLM: Mixed-precision LLM quantization with algorithm-system co-design.
\newblock In \emph{OpenReview}, 2025.

\bibitem{wei2025mcmoe}
W.~Wei et al.
\newblock Mixture compressor for Mixture-of-Experts LLMs gains more.
\newblock In \emph{ICLR}, 2025.

\bibitem{huang2025moequant}
M.~Huang et al.
\newblock MoEQuant: Enhancing quantization for Mixture-of-Experts large language models via expert-balanced sampling and affinity guidance.
\newblock \emph{arXiv preprint arXiv:2505.03804}, 2025.

\bibitem{xie2024quantmoebench}
Y.~Xie et al.
\newblock Examining post-training quantization for Mixture-of-Experts: A benchmark.
\newblock In \emph{NeurIPS Workshop on Efficient Natural Language and Speech Processing}, 2024.

\bibitem{bondarenko2023quantizable}
Y.~Bondarenko, M.~Nagel, and T.~Blankevoort.
\newblock Quantizable transformers: Removing outliers by helping attention heads do nothing.
\newblock In \emph{EMNLP}, 2023.

\bibitem{mlx2023}
Apple.
\newblock MLX: An array framework for Apple silicon.
\newblock \url{https://github.com/ml-explore/mlx}, 2023.

\bibitem{qwen3.5}
Qwen Team.
\newblock Qwen3.5 technical report.
\newblock \url{https://qwenlm.github.io/blog/qwen3.5/}, 2026.

\bibitem{xiao2023smoothquant}
G.~Xiao, J.~Lin, M.~Seznec, H.~Wu, J.~Demouth, and S.~Han.
\newblock SmoothQuant: Accurate and efficient post-training quantization for large language models.
\newblock In \emph{ICML}, 2023.

\bibitem{badri2024hqq}
M.~Badri and A.~Shaji.
\newblock Half-quadratic quantization of large machine learning models.
\newblock \url{https://github.com/mobiusml/hqq}, 2024.

\bibitem{huang2025slimllm}
W.~Huang et al.
\newblock SliM-LLM: Salience-driven mixed-precision quantization for large language models.
\newblock In \emph{ICML}, 2025.

\bibitem{shang2024cherryq}
Y.~Shang et al.
\newblock Cherry on Top: Parameter heterogeneity and quantization in large language models.
\newblock In \emph{NeurIPS}, 2024.

\bibitem{park2025hestia}
J.~Park et al.
\newblock HESTIA: A Hessian-guided differentiable quantization-aware training framework for extremely low-bit LLMs.
\newblock \emph{arXiv preprint arXiv:2601.20745}, 2026.

\bibitem{xu2025qwen3quant}
Z.~Xu et al.
\newblock An empirical study of Qwen3 quantization.
\newblock \emph{arXiv preprint arXiv:2505.02214}, 2025.

\bibitem{llama4}
Meta AI.
\newblock The Llama 4 herd of models.
\newblock \url{https://ai.meta.com/blog/llama-4-multimodal-intelligence/}, 2025.

\bibitem{eval-harness}
L.~Gao, J.~Tow, B.~Abbasi, S.~Biderman, S.~Black, A.~DiPofi, C.~Foster, L.~Golding, J.~Hsu, A.~Le~Noach, H.~Li, K.~McDonell, N.~Muennighoff, C.~Ociepa, J.~Phang, L.~Reynolds, H.~Schoelkopf, A.~Skowron, L.~Sutawika, E.~Tang, A.~Thite, B.~Wang, K.~Wang, and A.~Zou.
\newblock A framework for few-shot language model evaluation.
\newblock \url{https://github.com/EleutherAI/lm-evaluation-harness}, 2024.

\end{thebibliography}

% ═══════════════════════════════════════════════════════════════════════════════
% Appendix
% ═══════════════════════════════════════════════════════════════════════════════

\appendix

\section{Gap Closure Summary}
\label{app:gap_closure}

\begin{table}[h]
\centering
\caption{Perplexity gap closure: fraction of the quality gap between uniform 4-bit and BF16 that MINT recovers at various budget points.}
\label{tab:gap_closure}
\begin{tabular}{llrrr}
\toprule
Model & Budget & MINT $\Delta$PPL & Uniform $\Delta$PPL & Gap Closed \\
\midrule
Qwen3-8B & 6.0\,GB & +3.2\% & +5.4\% & 41\% \\
\midrule
Qwen3-30B & 16.3\,GB & +2.3\% & +10.3\% & 78\% \\
 & 17.4\,GB & +1.5\% & +10.3\% & 86\% \\
 & 19.0\,GB & +0.6\% & +10.3\% & 94\% \\
\midrule
GLM-4.7-Flash & 15.8\,GB & +5.8\% & +31.6\% & 82\% \\
\midrule
Llama-4-Scout & 58.0\,GB & $-$2.5\%$^*$ & --- & --- \\
Llama-4-Scout & 163.2\,GB & $-$6.8\%$^*$ & --- & --- \\
\bottomrule
\end{tabular}
\end{table}

$^*$Llama-4-Scout's BF16 model (${\sim}$203\,GB) exceeds available memory; $\Delta$ computed vs uniform 4-bit (PPL 7.899). Gap closure is not computable without BF16 baseline.

On Qwen3-30B-A3B, MINT closes 78--94\% of the perplexity gap depending on the budget, with the 19\,GB budget achieving near-BF16 quality at 33\% of full size.

\section{Comparison with Threshold-Based Allocation (v1)}
\label{app:v1_comparison}

MINT's MCKP formulation addresses specific limitations identified in the threshold-based v1 approach:

\begin{table}[h]
\centering
\small
\caption{Systematic comparison of threshold-based (v1) and optimization-based (MINT) design decisions.}
\label{tab:v1v2}
\begin{tabular}{p{3cm}p{4.5cm}p{4.5cm}}
\toprule
Aspect & Threshold-based (v1) & MINT (MCKP) \\
\midrule
Objective & Weighted sum + thresholds & Constrained optimization (Eq.~\ref{eq:objective}) \\
Error metric & Single-point 4-bit NRMSE & Multi-point RD curve (8 configs) \\
Group size & Fixed hyperparameter & Per-tensor variable (85\% chose g32) \\
Protection & Binary hard-coded rules & Soft priors in objective (Table~\ref{tab:priors}) \\
Safety floor & None (allows SQNR $<$ 5\,dB) & SQNR veto at 9\,dB (empirically validated) \\
Budget & No user control (fixed output) & User-specified (GB, avg bits, or ratio) \\
16-bit alloc & 5.6\% of params & 0\% (wasteful, redirected to g32 overhead) \\
2-bit alloc & 4.0\% of params & 0\% (blocked by SQNR floor) \\
Quality predict & Not possible & Fitted curve (Eq.~\ref{eq:ppl_curve}) \\
\bottomrule
\end{tabular}
\end{table}

\section{Exploratory Sensitivity Features}
\label{app:features}

In addition to the rate-distortion curves and SQNR values used by the allocator, MINT computes several per-tensor sensitivity features during analysis. These features were investigated as candidates for a per-tensor importance weight in the allocator objective, but the current system uses only raw NRMSE as the loss function. We document them here as they may be useful for future extensions that incorporate learned importance weighting.

\subsection{Spectral Features}

We compute three scale-invariant spectral features from the singular values $\sigma_1 \geq \sigma_2 \geq \cdots \geq \sigma_r$ (via randomized SVD with rank $k = 256$):

\paragraph{Stable rank.} The effective dimensionality of $\mathbf{W}$:
\begin{equation}
    r_s(\mathbf{W}) = \frac{\|\mathbf{W}\|_F^2}{\|\mathbf{W}\|_2^2} = \frac{\sum_{i=1}^{r} \sigma_i^2}{\sigma_1^2}
\end{equation}
Lower values indicate that fewer directions carry most information, making quantization more likely to corrupt critical components. Note that $\|\mathbf{W}\|_F^2$ is computed exactly from the weight tensor (not from the truncated SVD), while $\sigma_1$ comes from the top singular value.

\paragraph{Spectral tail mass.} The fraction of energy outside the top $k/10$ singular values:
\begin{equation}
    \tau(\mathbf{W}) = 1 - \frac{\sum_{i=1}^{\lfloor k/10 \rfloor} \sigma_i^2}{\|\mathbf{W}\|_F^2}
\end{equation}
The denominator uses the exact Frobenius norm, not the truncated singular value sum, ensuring the metric is well-defined regardless of the SVD rank parameter.

\paragraph{Approximate log spectral spread.} Based on the top-$k$ truncated SVD (not the true condition number):
\begin{equation}
    \kappa_k(\mathbf{W}) = \min\!\left(10, \, \log_{10}\!\left(\frac{\sigma_1}{\sigma_k + \varepsilon}\right)\right)
\end{equation}
Since only $k=256$ singular values are computed, this measures the spectral spread within the top-$k$ subspace rather than the true condition number.

\subsection{Per-Group Kurtosis Features}

We reshape $\mathbf{W}$ into $K = \lceil mn / g \rceil$ groups of size $g = 128$ and compute the excess kurtosis $\kappa_j$ for each group $j$:
\begin{equation}
    \kappa_j = \frac{1}{g} \sum_{i=1}^{g} \left(\frac{w_{j,i} - \bar{w}_j}{s_j}\right)^4 - 3
\end{equation}

Four features are extracted from the distribution $\{\kappa_1, \ldots, \kappa_K\}$:
\begin{align}
    f_{\text{kurt90}} &= P_{90}(\kappa_1, \ldots, \kappa_K) & \text{(90th percentile)} \\
    f_{\text{spread}} &= P_{99} - P_{50} & \text{(tail heaviness spread)} \\
    f_{\text{outlier}} &= \frac{1}{K} \sum_{j=1}^{K} \mathbf{1}[\exists\, w_{j,i} : |w_{j,i} - \bar{w}_j| > 3s_j] & \text{(outlier group fraction)} \\
    f_{\text{maxmed}} &= \max(\kappa_j) - \text{median}(\kappa_j) & \text{(max-to-median gap)}
\end{align}
Per-group kurtosis aligns the statistical metric with the quantization granularity, addressing the critique that global kurtosis operates at the wrong level~\cite{zhao2025kurtail}.

\subsection{Norm-Guided Output Noise Amplification}

Each linear layer is preceded by a normalization layer with a learned scale parameter $\boldsymbol{\gamma} \in \mathbb{R}^n$. We sample probes from $\mathbf{x}_j \sim \mathcal{N}(\mathbf{0}, \text{diag}(\boldsymbol{\gamma}^2))$, encoding the model's own channel importance without calibration data. Using the actual RTN quantization residual $\Delta \mathbf{W} = Q(\mathbf{W}; 4, 128) - \mathbf{W}$, the output noise amplification is:
\begin{equation}
    f_{\text{out}} = \frac{\|\Delta \mathbf{W} \cdot \mathbf{X}\|_F}{\|\mathbf{W} \cdot \mathbf{X}\|_F + \varepsilon}
\end{equation}
averaged over $p = 32$ probe vectors. We note that $f_{\text{out}}$ uses a fixed 4-bit residual as the perturbation; unlike single-point allocation methods, this feature is auxiliary---the allocator's actual loss function is the full multi-point RD curve, so the single-point circularity critique does not apply to the allocation decision. If a preceding normalization layer is unavailable, we fall back to isotropic probes $\mathbf{x} \sim \mathcal{N}(\mathbf{0}, \mathbf{I})$.

\section{Reproduction}
\label{app:reproduction}

All code is available at [anonymized]. Experiments use:
\begin{itemize}
    \item \texttt{python 3.12.0}, \texttt{mlx 0.30.3}, \texttt{mlx\_lm 0.30.4}, \texttt{torch 2.6.0}, \texttt{scipy}
    \item Perplexity: custom evaluator, WikiText-2 test split, seq\_len=2048, seed=42
    \item Quantization simulation: group-wise RTN
    \item MCKP solver: greedy efficiency ordering (default)
    \item SQNR floor: 9\,dB (empirically validated; blocks 2-bit, permits 3-bit)
\end{itemize}

\end{document}
